\section{Self-consistent fields}
\label{sec:The-iterative-procedure}

The effect of the backreaction of the scalar field to the background
electromagnetic field is studied by an iterative procedure. 
For each value of $\lambda$, at each iteration $\kappa$, one computes the following
\begin{subequations}
   \label{eq:iterative-procedure}
\begin{align}
    \frac{d {^2}\phi_{n, \lambda}^{\kappa}}{d {z^2}} {}  &= 
    -\left[ (\omega_{n, \lambda}^\kappa - \varepsilon A^{\kappa -1 }_{0, \lambda} )^2 
     - a^2 m^2\right] \phi_{n, \lambda}^{\kappa}, 
    \label{eq:mode-equation-iterative}
    \\
\begin{split}
        \rho_\lambda^{\kappa} &= \lim_{\tau \to 0} \Big[\sum_{n>0}^{} (\omega_{n,
        \lambda}^\kappa - \varepsilon A_{0, \lambda}^{\kappa-1} )\lvert \phi_{n,
        \lambda}^\kappa\rvert^2 e^{-i\omega^\kappa_{n, \lambda} \tau}\\
        &+\sum_{n<0}^{} (\omega^{\kappa}_{n, \lambda} - \varepsilon
        A_{0, \lambda}^{\kappa-1}) \lvert \phi^\kappa_{n, \lambda}\rvert^2
    e^{i\omega^{\kappa}_{n, \lambda} \tau}\Big] \\
    &+
    \frac{\varepsilon^2}{\pi}    A_{0, \lambda}^{\kappa-1} ,
\end{split}
\label{eq:vacuum-polarization-iterative}\\
\begin{split}
        A^{\kappa }_{0, \lambda} &= c A_{0, \lambda}^{\kappa-1} \\&+ \left(1-c
        \right)\left[ -\lambda\left(z - \frac{1}{2} \right) -
            \int_{\frac{1}{2}}^{z} \int_{0}^{\chi} \rho_\lambda^{\kappa}(\xi) d\xi d\chi
             \right] .
\end{split}
        \label{eq:A0-general-iterative}
\end{align}
\end{subequations}
This procedure has two main differences from the one used in \cite{ambjorn_properties_1983}: the expression for the vacuum polarization, and the $c<1$ prescription in Eq.~\eqref{eq:A0-general-iterative}.
The details on how $\rho$ is calculated are explained in Appendix~\ref{sec:Computing-the-vacuum-polarization}. Initial conditions of this procedure are discussed below.

Noticing how Eq.~\eqref{eq:A0-general-iterative} acts as an update function  on $A_{0,\lambda}^{\kappa-1}$, we write this scheme as an infinite dimensional fixed point problem. This is, for each $\lambda$, we look for the fixed point $A_{0, \lambda}$ obeying
\begin{align}
    A_{0, \lambda} = f_\lambda(A_{0, \lambda}),
    \label{eq:fixed-point-problem}
\end{align}
with $f_\lambda$ given by the composition of steps \eqref{eq:mode-equation-iterative}, \eqref{eq:vacuum-polarization-iterative} and \eqref{eq:A0-general-iterative}. 
The fixed points of $f_\lambda$, denominated ``self-consistent solutions", are looked for by studying the limit $A_{0,\lambda}=\lim_{\kappa\to \infty}A_{0,\lambda}^{\kappa}$, where 
\begin{align}
    A_{0,\lambda}^{\kappa + 1} = f_\lambda(A_{0,\lambda}^{\kappa}) =
    f_\lambda^{\kappa}\left(A_{0, \lambda}^{0} \right), \textrm{with } f_\lambda^\kappa =
    \underbrace{f_\lambda \circ ... \circ f_\lambda}_{k \textrm{ times}}.
    \label{eq:self-consistent-limit}
\end{align}
If the limit $A_{0, \lambda}$ exists, it is given the name
self-consistent potential, as well as the corresponding self-consistent vacuum
polarization $\rho_{\lambda}$, field modes $\phi_{n,
\lambda}$ and mode energies $\omega_{n, \lambda}$.

For low enough $\lambda\ll \lambda_c$, the initial value of the sequence $A_{0,
\lambda}^0$  can be set to the external field approximation $A_{0,
\lambda}^0(z) = -\lambda \left(z- \frac{1}{2}\right).$ This is, however,  no longer a valid initial potential for stronger background fields,  as it leads to non-convergences in the sequence $ A_{0, \lambda}^{\kappa}$. This is discussed in the Appendix ??(to be added). For higher $\lambda$-values,
$A_{0, \lambda}^{0}$ is instead set by a stored previously calculated self-consistent
potential $A_{0, \lambda'}$ with $\lambda = \lambda' + \Delta
\lambda$, where $\Delta \lambda > 0$. We use this scheme to probe the $\lambda$-parameter space.



Finally, in Eq.~\eqref{eq:A0-general-iterative}, $c\in [0, 1]$ is introduced in order to ``relax" the procedure: For
$c=0$, the procedure is precisely the one used in~\cite[]{ambjorn_properties_1983}, and it is expected to quickly converge. For
$c$ close to 1, updates will be small corrections to the previous potential.
%This avoids certain numerical instabilities appearing in the strong-field regime $\lambda \gtrsim \lambda_c$. 
Stronger background field strengths increase the instability of the coupled system, thus increasing the relevance of slower converge schemes. 
Considering this, we adopt a scheme in which, as $\lambda$ increases, $c$ increases and simultaneously $\Delta \lambda$ decreases.

%The mode subindices have been removed to avoid notation clutter.
%The iterative procedure is schematized in Fig.~\ref{fig:iterative-procedure},
%and it studies $\lim_{\kappa \to \infty} A_\lambda^\kappa $ for each $\lambda$.
%If the limit exists, it defines the self-consistent potential
%$A^{\infty}_\lambda$. Additionally, every other relevant quantity
%$\rho^\infty_\lambda, \omega_\lambda^{\infty}, \text{ and } \phi_\lambda
%^{\infty}$ also acquire the adjective ``self-consistent".

\begin{figure}[t]
\begin{center}
\begin{tikzpicture}[
    scale=0.87,
    box/.style = {draw, rectangle, minimum width=2cm, minimum height=1cm, fill={rgb:red,1;green,2;blue,2}, fill opacity=0.2, text opacity = 1},
    arrow/.style = {->, thick, >=Stealth},
    curved arrow/.style = {->, thick, >=Stealth, bend left}
]
% Nodes
\node[box] (A0) at (0, 0) {$A_{\lambda}^{0}$};
\node[box] (phi) at (3, 0) {$\omega_{n, \lambda}^\kappa, \phi_{n, \lambda}^\kappa$};
\node[box] (rho) at (8.25, 3) {$\rho_\lambda^\kappa$};
\node[box] (tildeA0) at (8.25, -3) {$A_\lambda^{\kappa}$};
\node (kappa) at (6.5, 0) {$\kappa \to \infty$};

% Arrows
\draw[arrow] (A0) -- (phi);

\draw[curved arrow, blue] (phi) to[out=50, in=135] (rho);
\draw[curved arrow, blue] (rho) to[out=55, in=135] (tildeA0);
\draw[curved arrow, blue] (tildeA0) to[out=45, in=130] (phi);

  %\draw[arrow] (phi) to[out=70, in=150, looseness=1] (rho);


% Central evolution arrow
\draw[->, thick, dashed, red] (8, 0) arc[start angle=360, end angle=0, radius=1.5cm];
%\draw[->, thick, dashed, red] (10, 0) arc[start angle=0, end angle=360, radius=3.5cm];

\end{tikzpicture}
\end{center}
\caption{Schematic of the iterative procedure.}%
\label{fig:iterative-procedure}
\end{figure}




%For a low $\lambda\ll \lambda_c$, the initial potential $A_\lambda^{0} = -\lambda
%\left(z-\frac{1}{2}\right)$ is given by the external field approximation. For
%higher values of $\lambda$, $A_\lambda^{0}$ is given by an already calculated
%self-consistent $A_{\lambda'}^{\infty}$, with $\lambda = \lambda' + \Delta
%\lambda$, for $\Delta \lambda > 0 $. 
